\chapter[Recursive Supercombinators]{Recursive\\Supercombinators}
%\vspace{2cm}
\vspace{1cm}

So far we have made no explicit mention of how our lambda-lifter should
handle recursive definitions. One way to do so is to translate all our recursive
definitions into non-recursive ones, using the fixpoint combinator \ml{Y}, as
described in Chapter~2. This is inefficient and slow for the following reasons:

\begin{numbered}
    \item There is no reason why the \textit{supercombinators} should not be explicitly
    recursive since, unlike lambda abstractions, they have names so they can
    refer to themselves. For example

    \begin{letalign}
        \ml{\$F x = \$G (\$F ($-$ x 1)) 0}
    \end{letalign}

    \item To make \ml{\$F} non-recursive using \ml{Y} would require an auxiliary definition,
    thus:

    \begin{letalign}
        \ml{\$F} &= \ml{Y \$F1} \\
        \ml{\$F1 F x} &= \ml{\$G (F ($-$ x 1)) 0}
    \end{letalign}

    Defining \ml{\$F} in this way will require more reductions than the explicitly
    recursive version, since the \ml{Y} has to be reduced.

    \item In Chapter~6 the translation into the ordinary lambda calculus of a \ml{letrec}
    involving mutual recursion was handled by first grouping the definitions
    into a tuple, and then making this definition non-recursive with \ml{Y}. Not
    only is it annoying to have to introduce tuples to handle mutual recursion
    of functions, but it is also very inefficient since the tuple has to be
    constructed and then taken apart.
\end{numbered}
We conclude that explicitly recursive definitions of supercombinators will
give a better performance. We now describe the techniques required to obtain
a set of mutually recursive supercombinator definitions without using \ml{Y}.

\section{Notation}

Since we want to treat recursion directly, we do not want it to be compiled into
applications of the \ml{Y} combinator. Hence we assume that the high-level
functional program is instead translated into the lambda notation augmented
with the simple let and letrec constructs, as was described in Chapter~3.

In passing we observe that the notation

\combinatorbox{
    \ml{\$S1 x y = B1} \\
    \indent \ml{\$S2 f = B2} \\
    \indent \ml{etc.}
}{
    \ml{E}
}

\noindent
is precisely equivalent to
\begin{letalign}
        letrec & \\
        & \ml{\$S1} = \tlb{x}\tlb{y}B1 \\
        & \ml{\$S2} = \tlb{f}B2 \\
        & etc. \\
        in     & \\
        & \ml{E}
\end{letalign}
so that the lambda-lifting process can be regarded as a source to source
transformation of the enriched lambda calculus.

\section{\ml{Let}s and \ml{letrec}s in Supercombinator Bodies}

Suppose we wanted to write a textual description for the graph

\plainbox{\centering
    \vspace{-0.5\baselineskip}
\begin{tikzpicture}
    \node (a) at (0,   0)  {\ml{@}};
    \node (b) at (0.5,  -1)  {\ml{@}};
    \node (c) at (-1, -2)  {\ml{@}};
    \node (f) at (-1.7,-3) {\ml{f}};
    \node (g) at (-0.3,-3) {};
    \node (three) at (1, -2) {\ml{3}};
    \draw (a) -- (b);
    \draw (a) -- (c);
    \draw (b) -- (c);
    \draw (b) -- (three);
    \draw (c) -- (f);
    \draw[-{Stealth[length=3mm,width=2mm]}] (c) -- (g.center) -- (1.8,-3) -- (1.8,-1) -- (b);
\end{tikzpicture}
    \vspace{-0.5\baselineskip}
}

\noindent
Whilst expressions such as \ml{(f (g a) b)} can describe trees, they cannot express
the sharing and cycles embodied in the above graph. One solution would be to
name the nodes (\ml{a}, \ml{b} and \ml{c}, say, working top to bottom) and express the
graph thus:
\begin{mlcoded}
    a = c b \\
    b = c 3 \\
    c = f b
\end{mlcoded}
We would also want to identify \ml{a} as being the root of the graph. But we have
just reinvented the \ml{letrec}! The graph can be described by the \ml{letrec} expression
\begin{mlcoded}
    \begin{tabular}{@{}l@{\;}l@{}}
        letrec & a = c b \\
        & b = c 3 \\
        & c = f b \\
        in a
    \end{tabular}
\end{mlcoded}
This gives us the idea that \ml{letrec} expressions can be regarded as \textit{textual
    descriptions of a cyclic graph}. Hence a \ml{letrec} in a supercombinator body can be
regarded as the description of a graphical portion of the supercombinator
body.

Up to now we have considered a supercombinator body to be a \textit{tree}, and
applying the supercombinator involves constructing a new instance of the
tree. Now we see that allowing \ml{letrec}s in a supercombinator body allows the
body to be a \textit{graph}, and applying such a supercombinator involves con-
structing a new instance of this graph. We say that such a supercombinator has
a \textit{graphical body}.

For example, consider the following supercombinator definition:
\begin{letalign}
        \$Y f = & letrec yf = f yf \\
        & in yf
\end{letalign}
This is a definition of the cyclic version of the familiar \ml{Y} combinator, whose
body is a graph. When \ml{\$Y} is applied, we make an instance of the graph,
substituting for occurrences of the formal parameter, \ml{f}. During the
instantiation we must be careful to preserve the cycles of the original graph.

A compiling implementation would compile code which would, when
executed, construct the graph with the appropriate substitutions made. The
way in which this is done is described in Chapter~18.

In a similar way we can allow supercombinator bodies to contain \ml{let}-expressions, regarding them as descriptions of (acyclic) graphs. This will
actually save us reductions, because we can now describe directly expressions
such as
\begin{mlcoded}
    let x = 3 in E
\end{mlcoded}
where we would previously have translated this to
\begin{letalign}
    (\tlb{x}E) 3
\end{letalign}
which requires a reduction to explicate.

To summarize, we see that

\begin{numbered}
    \item it is quite easy to extend supercombinators to allow them to have bodies
    which are general graphs, rather than being restricted to trees;
    \item graphical supercombinator bodies can easily be described using a \ml{letrec}
    (or a \ml{let} in the case of acyclic bodies);
    \item to instantiate a \ml{letrec} (or \ml{let}), we simply construct the graph described by
    the \ml{letrec} (or \ml{let});
    \item using graphical bodies can save us reductions.
\end{numbered}
We now discuss how to transform recursive programs into supercombinators
with graphical bodies.

\section{Lambda-lifting in the Presence of \ml{letrec}s}

Our lambda-lifting algorithm will work as before, lifting the lambda
abstractions to the top level. No special note need be taken of \ml{letrec}s; they can
be treated just like any other expression. In particular, lambda-lifting still
applies only to lambda abstractions, not to \ml{letrec}s as well. Some lambda
abstractions will have \ml{letrec}s in their bodies, which will give rise to super-
combinators with graphical bodies.

The question arises, however, of what lexical level-number to assign to
variables bound in a \ml{letrec}.

The variables bound in a \ml{letrec} will be instantiated when the \textit{immediately
    (textually) enclosing lambda abstraction} is applied to an argument, since that
is when we construct the instance of the \ml{letrec}, substituting for all the free
variables. Hence the variables bound in a \ml{letrec} should be given the lexical
level-number of the immediately enclosing lambda abstraction.

What if there is no enclosing lambda abstraction? In this case the natural
level-number for such variables should be 0. But this gives us a hint, since 0 is
the level-number we assign to constants and supercombinators. If there is no
enclosing lambda abstraction, then the definition bodies of the \ml{letrec} can have
no free variables (other than the variables defined in the \ml{letrec}); in other words,
they are combinators. All that is needed to turn them into supercombinators
is to lambda-lift them to remove any inner lambdas. Notice that the variables
bound in such level 0 \ml{letrec}s will not be taken out as free variables because
constants (level 0) are not taken out.

Suppose we have to lambda-lift this program, which computes the infinite
list of 1s.

\combinatorbox{}{
    \ml{letrec x = CONS 1 x} \\
    \indent \ml{in x}
}

\noindent
The \ml{letrec} is at level 0, and there are no lambda abstractions, so \ml{x} is a
supercombinator already, and we get

\combinatorbox{
    \ml{\$x = CONS 1 \$x}
}{
    \ml{\$x}
}

As an example of a recursive function, consider the factorial function:

\combinatorbox[0.7\textwidth]{}{
    \vspace{-0.5\baselineskip}
    \begin{mlcoded}
        letrec fac = \tlb{n}IF (= n 0) 1 ($\ast$ n (fac ($-$ n 1))) \\
        in fac 4
    \end{mlcoded}
}

\noindent
The \ml{letrec} is at level 0, and there are no lambda abstractions inside the body of
the \ml{\tl n} abstraction. Hence, \ml{fac} is already a supercombinator and we get

\combinatorbox[0.7\textwidth]{
    \begin{mlcoded}
        \$fac n = IF (= n 0) 1 ($\ast$ n (\$fac ($-$ n 1))) \\
        \$Prog = \$fac 4
    \end{mlcoded}
        \vspace{-0.5\baselineskip}
}{
    \ml{\$Prog}
}

\section{Generating Supercombinators with Graphical Bodies}

So far none of our supercombinators has had a graphical body. This occurs
when a \ml{letrec} has some free variables. Consider, for example, the program

\combinatorbox[0.7\textwidth]{}{
        \vspace{-0.5\baselineskip}
    \begin{letalign}
            let & \\
            & Inf = \tlb{v}(letrec vs = CONS v vs in vs) \\
            in  & \\
            & Inf 4
    \end{letalign}
}

\noindent
\ml{(Inf v)} returns the infinite list of \ml{v}s. Again, \ml{Inf} is at level 0 and contains no inner
lambda abstractions, so it is already a supercombinator, and we get

\combinatorbox[0.7\textwidth]{
    \begin{mlcoded}
        \$Inf v = letrec vs = CONS v vs in vs \\
        \$Prog = \$Inf 4
    \end{mlcoded}
    \vspace{-0.5\baselineskip}
}{
    \ml{\$Prog}
}

\noindent
Notice that the graphical body of the supercombinator preserves the (finite)
cyclic representation of the (infinite) data structure.

\section{An Example}

We shall now work through an example to show the lambda-lifting algorithm
in action. Here is a Miranda program to sum the first 100 integers. It is written
in a slightly odd style in order to demonstrate various aspects of the algorithm:

\combinatorbox[0.9\textwidth]{
    \begin{mlcoded}
        \begin{tabular}{@{}l@{ }l@{ }l}
            sumInts m &= sum (count 1) & \\
            & \phantom{=} where & \\
            & \qquad\qquad count n  &= [\,], \qquad\qquad\qquad\qquad n>m \\
            &  &= n : count (n+1) \\
            sum [\,]    & = 0  &\\
            sum (n:ns)  & = n + sum ns &
        \end{tabular}
    \end{mlcoded}
%    \vspace{-0.5\baselineskip}
}{
    \ml{sumInts 100}
}

\noindent
Translating this into the enriched lambda calculus gives

\combinatorbox[0.9\textwidth]{
}{
    \vspace{-0.5\baselineskip}
    \begin{mlcoded}
        \begin{tabular}{l@{}l@{}}
            letrec  \\
            \quad sumInts  \\
            \quad = \tlb{m}letrec  \\
            \qquad\qquad\qquad count  = \tlb{n} &IF ($>$ n m) \\
            \quad\quad &NIL \\
            \quad\quad &(CONS n (count (+ n 1))) \\
            \qquad\qquad\ in  \\
            \qquad\qquad\qquad sum (count &1)  \\
            \quad sum  = \tlb{ns}IF (= ns &NIL) 0 (+ (HEAD ns) (sum (TAIL ns))) \\
            in  \\
            \quad sumInts 100
        \end{tabular}
    \end{mlcoded}
    \vspace{-0.5\baselineskip}
}

\noindent
(Note: this is not exactly the translation that will be produced by the pattern-matching compiler described in Chapter~5, but it is a correct translation, and will suffice for present purposes.) The variables \ml{sumInts} and \ml{sum} are defined at level~0, but \ml{sumInts} has an inner lambda abstraction. This \ml{\tl n} abstraction has the free variables \ml{m} and \ml{count}. We lift them out to generate a supercombinator, which we arbitrarily name \ml{\$count}, thus

\combinatorbox[0.9\textwidth]{
    \begin{mlcoded}
        \$count count m n = IF ($>$ n m) NIL (CONS n (count (+ n 1)))
    \end{mlcoded}
    \vspace{-0.5\baselineskip}
}{
    \vspace{-0.5\baselineskip}
    \begin{mlcoded}
        \begin{tabular}{@{}l@{ }l}
            letrec & \\
            \qquad sumInts & \\
            \qquad = \tlb{m} & letrec  \\
            & \qquad count = \$count count m \\
            & in  \\
            &\qquad sum (count 1) \\
            \qquad sum  = &\tlb{ns}IF (= ns NIL) 0 (+ (HEAD ns) (sum (TAIL ns))) \\
            in  \\
            \qquad sumInts &100
        \end{tabular}
    \end{mlcoded}
    \vspace{-0.5\baselineskip}
}

Now \ml{sumInts} and \ml{sum} have no inner abstractions, and they are at the top
level, so they are supercombinators. Lifting them directly, and adding the
final \ml{\$Prog} supercombinator, gives

\combinatorbox[0.9\textwidth]{
    \begin{mlcoded}
        \$count count m n = IF ($>$ n m) NIL (CONS n (count (+ n 1))) \\
        \$sum ns = IF (= ns NIL) 0 (+ (HEAD ns) (\$sum (TAIL ns))) \\
        \$sumInts m = letrec count = \$count count m \\
        \phantom{\$sumInts m = }in \$sum (count 1) \\
        \$Prog = \$sumInts 100
    \end{mlcoded}
    \vspace{-0.5\baselineskip}
}{
    \ml{\$Prog}
}

\noindent
We are done.

\section{Alternative Approaches}

The technique described earlier is not the only way of lambda-lifting recursive
functions. For example, Johnsson [1985] describes an algorithm which
constructs graphical supercombinator bodies for data structures, but not for
functions.

Briefly, his technique works like this. Suppose we have a program with a
recursive function \ml{f} containing a free variable \ml{v}:

\combinatorbox{}{
    \vspace{-0.5\baselineskip}
    \begin{mlcoded}
        (\ldots \\
        \phantom{XX} letrec f = \tlb{x}(\ldots f\ldots v\ldots) \\
        \phantom{XX} in (\ldots f\ldots) \\
        \ldots)
    \end{mlcoded}
    \vspace{-0.5\baselineskip}
}

\noindent
We generate a recursive supercombinator \ml{\$f} from \ml{f} by abstracting the free
variables (just \ml{v} in this case) \textit{but not} \ml{f} \textit{itself}. Instead, all uses of \ml{f} are replaced
with \ml{(\$f v)}, including those in the body of \ml{\$f} itself. This yields

\combinatorbox{
    \begin{mlcoded}
        \$f v x = \ldots(\$f v)\ldots v\ldots
    \end{mlcoded}
    \vspace{-0.5\baselineskip}
}{
    \vspace{-0.5\baselineskip}
    \begin{mlcoded}
        (\ldots \\
        \phantom{XX} (\ldots(\$f v)\ldots) \\
        \ldots)
    \end{mlcoded}
    \vspace{-0.5\baselineskip}
}

To illustrate the method we will recompile the \ml{sumInts} example given
earlier. Recall that we begin with the program

\combinatorbox[0.9\textwidth]{
}{
    \vspace{-0.5\baselineskip}
    \begin{mlcoded}
        \begin{tabular}{l@{}l@{}}
            letrec  \\
            \quad sumInts  \\
            \quad = \tlb{m}letrec  \\
            \qquad\qquad\qquad count  = \tlb{n} &IF ($>$ n m) \\
            \quad\quad &NIL \\
            \quad\quad &(CONS n (count (+ n 1))) \\
            \qquad\qquad\ in  \\
            \qquad\qquad\qquad sum (count &1)  \\
            \quad sum  = \tlb{ns}IF (= ns &NIL) 0 (+ (HEAD ns) (sum (TAIL ns))) \\
            in  \\
            \quad sumInts 100
        \end{tabular}
    \end{mlcoded}
    \vspace{-0.5\baselineskip}
}

First we lambda-lift the \tl n abstraction, abstracting out the free variable \ml{m},
\textit{but not} \ml{count}. Instead, we replace all calls to \ml{count} with \ml{(\$count m)}, which gives

\combinatorbox[0.9\textwidth]{
        \ml{\$count m n = IF ($>$ n m) NIL (CONS n (\$count m (+ n 1)))}
}{
    \vspace{-0.5\baselineskip}
    \begin{letalign}
            letrec  &\\
            & sumInts  \\
            & = \tlb{m}sum (\$count m 1)  \\
            & sum  = \tlb{ns}IF (= ns NIL) 0 (+ (HEAD ns) (sum (TAIL ns))) \\
            in & \\
            & sumInts 100
    \end{letalign}
}

There were two calls to \ml{count}, one in the body of the \tl n abstraction and one in
the definition of \ml{sumInts}, both of which were replaced with \ml{(\$count m)}. Notice
that this substitution could equally well be carried out using a let-definition to
bind \ml{count} to \ml{(\$count m)}.

Now \ml{sumInts} and \ml{sum} are supercombinators, so we lift them out to give

\combinatorbox[0.9\textwidth]{
    \begin{mlcoded}
        \$count m n = IF ($>$ n m) NIL (CONS n (\$count m (+ n 1))) \\
        \$sum ns = IF (= ns NIL) 0 (+ (HEAD ns) (\$sum (TAIL ns))) \\
        \$sumInts m = \$sum (\$count m 1) \\
        \$Prog = \$sumInts 100
    \end{mlcoded}
    \vspace{-0.5\baselineskip}
}{
    \ml{\$Prog}
}

\noindent
Notice that (unlike our previous method) no supercombinator has a graphical
body; all the recursion is handled by direct recursion of supercombinators.
However, it turns out that cyclic data structures have to be treated in a
different way, and do require supercombinators with graphical bodies.

The new method has one major advantage. In our previous approach the
recursive call to \ml{count} in the \ml{\$count} supercombinator was made to a function
passed in as a parameter (called \ml{count}). In contrast, the new method makes
the recursive call directly to the supercombinator \ml{\$count}. This means that the
compiler can see which function is being called, and this information can
make the compiled code considerably more efficient (see Chapter~20).

On the other hand, the \ml{\$count} supercombinator generated by the new
method is larger than that generated by the previous method. It contains an
extra application node \ml{(\$count m)}, and a new instance of this application node
will be constructed on every application of \ml{\$count}, which will consume more
store.

In the case of mutually recursive functions, it turns out that each function
needs to be passed the free variables of all the other functions in the mutually
recursive set, as well as its own. This involves doing the sort of dependency
analysis described in Section~6.2.8. Furthermore, as mentioned above, data
structures and functions must be treated in different ways by the new method,
which makes the compiler more complicated.

The trade-off between the two techniques is not yet clear.

\section{Compile-time Simplifications}

Once lambda-lifting has been completed there are some simple optimizations
that further improve the lambda-lifted program. These take the form of
compile-time simplifications of the program.

\subsection{Compile-time Reductions}

It may be advantageous to perform certain reductions at compile-time. For
example, consider the definitions
\begin{mlcoded}
    \$F x y = + (\$G y) x \\
    \$G p = $\ast$ p p
\end{mlcoded}
The \ml{(\$G y)} in the body of \ml{\$F} is a redex which will be created every time \ml{\$F} is
applied. We could, however, reduce it at compile-time, giving
\begin{mlcoded}
    \$F x y = + ($\ast$ y y) x
\end{mlcoded}
thus performing the \ml{\$G} reduction once and for all at compile-time. This
process is directly analogous to expanding out the code for a procedure call
in-line, a common optimization in conventional compilers. In order to
preserve sharing we should replace the redex with a let-expression:
\begin{mlcoded}
    \$G E $\rightarrow$ let p = E in $\ast$ p p
\end{mlcoded}
Sometimes we can evaluate subexpressions completely, as in the definition
\begin{mlcoded}
    \$H x = + x ($\ast$ 3 4)
\end{mlcoded}
where we can safely evaluate the \ml{($\ast$ 3 4)} once and for all at compile-time.
This is called constant folding in conventional compiler technology.

The question of exactly which redexes to reduce is not completely straight-forward, especially in the case of recursive functions, because indiscriminate use of the technique may cause the code size to increase significantly. The decision is not clear-cut, because it depends on the relative priorities of speed and code size. Hudak and Kranz [1984] give an interesting discussion of a particularly thorough-going use of compile-time reduction.

\subsection{Common Subexpression Elimination}

Sometimes (for clarity) the programmer may write an expression such as
\begin{mlcoded}
    $\ast$ (\$F x) (\$F x)
\end{mlcoded}
Rather than compute \ml{(\$F x)} twice, we can replace the expression with
\begin{mlcoded}
    let fx = \$F x in $\ast$ fx fx
\end{mlcoded}

Identifying common subexpressions may be done by a hashing algorithm
which checks to see if an expression already exists before building a new one.
This simplification seems always to be beneficial, but see Chapter~23 for a
warning about some possible drawbacks.

\subsection{Eliminating Redundant \ml{let}s}

Sometimes \ml{let}s of the form
\begin{mlcoded}
    let x = y in E
\end{mlcoded}
arise, in which the right-hand side of the definition is a single variable. These
can safely be eliminated by replacing occurrences of \ml{x} by \ml{y} in \ml{E}.

It is also quite common to encounter code in which a variable defined in a \ml{let}
is used only once in its scope. For example, consider the supercombinator
\begin{mlcoded}
    \begin{tabular}{l@{ }l@{ }l}
        \$sumSq x y = &let & xsq  = $\ast$ x x \\
        & & ysq = $\ast$ y y \\
        & in & \\
        & &+ xsq ysq
    \end{tabular}
\end{mlcoded}
In this case we may as well substitute the right-hand side of the definition for
the (single) occurrence of the variable, giving
\begin{mlcoded}
    \$sumSq x y = + ($\ast$ x x) ($\ast$ y y)
\end{mlcoded}
This is simpler and, it turns out, slightly more efficient. It may be achieved
simply by accumulating information on the number of textually distinct
occurrences of each variable in the body of a \ml{let}.

 \section*{References}

\begin{references}
    \item Hudak, P., and Kranz, D. 1984. A combinator based compiler for a functional
    language. In \textit{Proceedings of the 11th ACM Symposium on Principles of
        Programming Languages}, pp. 122--32. January.

    \item Johnsson, T. 1985. Lambda-lifting -- transforming programs to recursive equations. In
    \textit{Conference on Functional Programming and Computer Architecture, Nancy}, pp.
    190--203. Jouannaud (editor). LNCS 201. Springer Verlag.
\end{references}